% Essential joint-log proof; explicit constants are expanded in Appendix B.
由于两拍载荷使用同一私有符号，需要计算其联合分布，才能检验完整日志约束。
给定攻击前完整历史，令
$S_h=P_k^-(h)+R$、$S'_h=AF_1(P_k^-(h))A^\top+Q+R$，
$u_h=AK_hv$，并以 $\phi_\Sigma$ 表示 $N(0,\Sigma)$ 的密度。
当前真实创新的未归一化发送密度为
$f_h(u)=\phi_{S_h}(u)[1-g(u)]$，其积分为 $p_h$。
给定当前真实发送测量，下一真实条件创新 $e$ 服从 $N(0,S'_h)$，
且该条件分布不随当前创新改变。
对有符号幅值 $t$，下一预测偏差为 $-tu_h$；
若下一时刻发送，补偿使接收创新恰为 $e$。
下一记录既可能是静默事件，也可能包含连续载荷。将两者共同表示为
$\{\bot\}\sqcup\R^n$ 上的概率核，并定义当前发送分支的联合核：
\begin{equation}
 \begin{aligned}
 \mathcal K_{h,d}(\{\bot\})&=\int\phi_{S'_h}(e)g(e+d)\,de,\\
 \mathcal K_{h,d}(de)&=\phi_{S'_h}(e)[1-g(e+d)]\,de,\\
 F_h(t)(dz,d\xi)&=f_h(z-tv)\,dz\,\mathcal K_{h,-tu_h}(d\xi).
 \end{aligned}
 \label{eq:main-joint-kernel}
\end{equation}
其中，$z$ 为当前接收创新，$F_h(t)$ 的总质量恒为 $p_h$。
式\eqref{eq:main-joint-kernel}的第一因子描述当前载荷的平移，第二因子描述预测反馈引起的下一事件概率变化。

记 $\|\cdot\|_{\rm var}$ 为有符号测度的变差范数，并令 $a_h=u_h^\top Zu_h$。
由 $\sup_z|\partial_{u_h}g(z)|\le\sqrt{a_h/\mathrm e}$、
$\sup_z|\partial_{u_h}^2g(z)|\le a_h$，
分别计入静默原子和连续密度的导数，得到
\begin{equation}
 \begin{aligned}
 \|F_h''(t)\|_{\rm var}
 &\le \|\partial_v^2 f_h\|_1
 +4\sqrt{a_h/\mathrm e}\|\partial_v f_h\|_1
 +2p_ha_h\\
 &\le B_h<\infty.
 \end{aligned}
 \label{eq:main-kernel-smooth}
\end{equation}
式\eqref{eq:main-kernel-smooth}中的三项依次来自当前密度的二阶导数、
两个因子的交叉导数及下一记录核的二阶导数。
$f_h$ 的高斯方向导数可积，平移保持其 $L^1$ 范数；
在共同支配测度 $dz\otimes(\delta_\bot+d\xi)$ 下，乘积核因而二阶连续 $L^1$ 可微。
固定模型和有限事件前缀使 $B_h$ 与变化的 $b,\varepsilon$ 无关；
附录\ref{app:full-kernel}给出其显式选择。
等概率正负符号使联合核的一阶扰动相消。
对 $Q_h^\eta=[F_h(\eta)+F_h(-\eta)]/2$ 使用中心 Taylor 公式，可得
\begin{equation}
 \begin{aligned}
 &Q_h^\eta-F_h(0)\\
 &\quad=\frac12\int_0^\eta(\eta-t)
 \bigl[F_h''(t)+F_h''(-t)\bigr]\,dt,\\
 &\|Q_h^\eta-F_h(0)\|_{\rm var}\le \tfrac12B_h\eta^2.
 \end{aligned}
 \label{eq:main-kernel-taylor}
\end{equation}
给定攻击前历史及 $(z,\xi)$，按
$m_{k+1}^-=A(m_k^-+K_hz)$，并在 $\xi\ne\bot$ 时令下一载荷为
$m_{k+1}^-+\xi$，即可重建全部事件、载荷和预测记录。
该重建映射不依赖 $t$ 或私有符号；当前静默分支则与无攻击运行相同。
对共同前史积分，利用 TV 在固定映射下的收缩性，以及等质量测度差的
TV 为其变差范数的一半，得到
$\Delta_{\rm full}\le C_\pm\eta^2$，
其中 $C_\pm=\frac14\sum_hw_hB_h\in(0,\infty)$。
